%----------------------------------------------------------------------------------------
% دستاوردهای منتخب
%----------------------------------------------------------------------------------------
\cvsection{دستاوردهای منتخب}

\begin{cventries}

  \cventry
    {هوش مصنوعی سلامت در مقیاس آمازون}
    {\textbf{\lr{Amazon Health AI}}}
    {۲۰۲۶}
    {}
    {\begin{cvitems}\item {مشارکت علمی کلیدی در توسعه و عرضه هوش مصنوعی سلامت آمازون؛ از راه‌اندازی دستیار سلامت مبتنی بر هوش مصنوعی عامل‌محور در \lr{One Medical} تا گسترش آن به اعضای \lr{Prime} و ارائه راهنمایی سلامت شخصی‌سازی‌شده به‌صورت شبانه‌روزی.}\end{cvitems}}

  \cventry
    {انتشار پژوهشی داوری‌شده}
    {\textbf{Peer-reviewed Publication}}
    {۲۰۲۶}
    {}
    {\begin{cvitems}\item {مشارکت در پژوهش و توسعه چارچوب‌ها و روش‌های ارزیابی برای عامل‌های هوش مصنوعی سلامت، با تمرکز بر قابلیت اتکا، ایمنی، عملکرد و سنجش رفتار عامل‌ها در سناریوهای واقعی.}\end{cvitems}}

  \cventry
    {ارزیابی و ایمنی هوش مصنوعی در محیط تولیدی}
    {\textbf{\lr{AI Safety \& Evaluation}}}
    {۲۰۲۴–۲۰۲۶}
    {}
    {\begin{cvitems}\item {طراحی و توسعه روش‌ها و چارچوب‌های ارزیابی برای سامانه‌های مبتنی بر \lr{LLM} و عامل‌های هوش مصنوعی، با تمرکز بر ایمنی، قابلیت اتکا، groundedness و شناسایی رفتارهای ناخواسته پیش از استقرار.}\end{cvitems}}

  \cventry
    {آموزش هوش مصنوعی و یادگیری ماشین}
    {\textbf{تدریس و توسعه دانشجو}}
    {۲۰۱۹–۲۰۲۱}
    {}
    {\begin{cvitems}\item {تدریس دروس پیشرفته یادگیری ماشین و یادگیری عمیق برای بیش از ۱۰۰ دانشجو در چهار نیم‌سال، هماهنگی بیش از ۵۰ گروه پروژه کارشناسی ارشد، و طراحی و اجرای کارگاه‌های عملی یادگیری عمیق و پردازش زبان طبیعی.}\end{cvitems}}

  \cventry
    {طراحی رویکرد آموزشی برای هوش مصنوعی نسل جدید}
    {\textbf{مبانی نظری \lr{→} پژوهش \lr{→} پروژه \lr{→} ارزیابی \lr{→} ایمنی و اخلاق}}
    {چشم‌انداز آموزشی}
    {}
    {\begin{cvitems}\item {تدوین رویکردی آموزشی برای پیوند مبانی نظری هوش مصنوعی با پژوهش تجربی، پروژه‌های واقعی، ارزیابی علمی، و مباحث ایمنی و اخلاق؛ با تمرکز بر فناوری‌هایی مانند \lr{LLM}، \lr{RAG} و \lr{Agentic AI}.}\end{cvitems}}

  \cventry
    {رهبری دانشگاهی}
    {\textbf{انجمن دانشجویان تحصیلات تکمیلی علوم کامپیوتر دانشگاه بوفالو}}
    {۲۰۲۰}
    {}
    {\begin{cvitems}\item {ریاست انجمن دانشجویان تحصیلات تکمیلی علوم کامپیوتر، ایجاد برنامه‌های ارتباطی میان دانشجویان، پژوهشگران و صنعت، و دریافت جایزه رهبری دانشجویان تحصیلات تکمیلی در سال ۲۰۲۰.}\end{cvitems}}

\end{cventries}
